
\documentclass[10pt,a4paper]{article}
\usepackage[utf8]{inputenc}
\usepackage[T1]{fontenc}
\usepackage[french]{babel}
\usepackage[left=0.6cm,right=0.6cm,top=0.4cm,bottom=0.4cm]{geometry}
\usepackage{amsmath}
\usepackage{amssymb}
\usepackage{array}
\usepackage{graphicx}
\usepackage{xcolor}
\usepackage{titlesec}
\usepackage{enumitem}
\usepackage{multicol}
\usepackage{pifont}
\usepackage{frcursive}
\usepackage{tikz}
\usetikzlibrary{calc,arrows.meta}
\usepackage{tcolorbox}
\tcbuselibrary{skins,breakable}
\usepackage[hidelinks]{hyperref}
\usepackage{fontawesome5}
\definecolor{calcblue}{HTML}{1C7ED6}
\definecolor{diffgreen}{HTML}{40C057}
\definecolor{diffblue}{HTML}{228BE6}
\definecolor{diffred}{HTML}{FA5252}
\definecolor{skyblue}{HTML}{339AF0}
\newcommand{\starfull}[1]{{\color{#1}\ding{72}}}
\newcommand{\starempty}{{\color{gray!45}\ding{73}}}
\newcommand{\Vj}[1]{\overrightarrow{\text{#1}}}
\newcommand{\VBlank}{\rule[0.15ex]{0.55cm}{0.35pt}}
\renewcommand{\familydefault}{\sfdefault}
\DeclareSymbolFont{sfnumbers}{OT1}{cmss}{m}{n}
\SetSymbolFont{sfnumbers}{bold}{OT1}{cmss}{bx}{n}
\DeclareMathSymbol{0}{\mathalpha}{sfnumbers}{`0}
\DeclareMathSymbol{1}{\mathalpha}{sfnumbers}{`1}
\DeclareMathSymbol{2}{\mathalpha}{sfnumbers}{`2}
\DeclareMathSymbol{3}{\mathalpha}{sfnumbers}{`3}
\DeclareMathSymbol{4}{\mathalpha}{sfnumbers}{`4}
\DeclareMathSymbol{5}{\mathalpha}{sfnumbers}{`5}
\DeclareMathSymbol{6}{\mathalpha}{sfnumbers}{`6}
\DeclareMathSymbol{7}{\mathalpha}{sfnumbers}{`7}
\DeclareMathSymbol{8}{\mathalpha}{sfnumbers}{`8}
\DeclareMathSymbol{9}{\mathalpha}{sfnumbers}{`9}
\setlength{\parindent}{0pt}
\graphicspath{{./}}
\newcommand*{\textcircledd}[1]{\begin{tikzpicture}[baseline=(letter.base)]\node[draw,circle,inner sep=1pt] (letter) {#1};\end{tikzpicture}}
\newcommand{\numeroCerne}[1]{\protect\textcircledd{\bfseries\scriptsize #1}}
\setlist[enumerate,1]{label=\numeroCerne{\arabic*}, leftmargin=*, labelsep=6pt, itemsep=3pt, topsep=2pt, partopsep=0pt, parsep=0pt}

\NewDocumentEnvironment{Exo}{m m m}{%
  \begin{tcolorbox}[
    enhanced,
    unbreakable,
    colback=white,
    colframe=white,
    boxrule=0pt,
    arc=4pt, outer arc=4pt, sharp corners=south,
    left=8pt, right=8pt, top=5pt, bottom=3pt,
    before skip=4pt, after skip=4pt,
  ]
  \noindent\tcbox[on line, colback=skyblue, colframe=skyblue, boxsep=1pt, left=4pt, right=4pt, top=1pt, bottom=1pt, arc=2pt]{\textcolor{white}{\bfseries Exercice~#1}}~#2\hfill\mbox{\color{skyblue}\footnotesize\itshape #3}\par\vspace{3pt}%
}{%
  \end{tcolorbox}
}

\newcommand{\Objectif}[2]{%
\vspace{4pt}
\noindent{\large\bfseries\color{calcblue}Objectif #1~:\\#2}\par
\vspace{2pt}
}

\NewDocumentCommand{\FeuilletTitre}{m}{%
  \begin{center}
  \begin{tcolorbox}[
    enhanced,
    width=0.78\linewidth,
    colback=white,
    colframe=black,
    boxrule=0.8pt,
    arc=8pt,
    sharp corners=downhill,
    left=10pt, right=10pt, top=8pt, bottom=6pt,
    halign=center,
    overlay={
      \node[fill=white, anchor=west, font=\normalsize, inner xsep=4pt]
        at ([xshift=14pt]frame.north west) {Feuille d'exercices n\textdegree{}2};
    },
  ]
  {\huge #1}
  \end{tcolorbox}
  \end{center}
}

\pagestyle{empty}
\begin{document}

\FeuilletTitre{Vecteurs}

\begin{multicols}{2}

\Objectif{1}{Vecteur $\Vj{MM'}$ associ\'e \`a la translation qui transforme $M$ en $M'$ --- direction, sens et norme}

\begin{Exo}{1}{\starfull{diffgreen}\starempty\starempty}{Repr\'esenter}
\`A partir de la figure ci-dessous,
\begin{enumerate}
\item donner les images des points $C$, $D$, $E$ dans la translation de vecteur $\Vj{AB}$~;
\item en d\'eduire trois vecteurs \'egaux au vecteur $\Vj{AB}$.
\end{enumerate}
\begin{center}
\begin{tikzpicture}[scale=0.38]
\draw[thin,gray!35] (-0.5,-0.5) grid (7.5,5.5);
\coordinate (A) at (1,4); \coordinate (B) at (4,5);
\coordinate (C) at (2,2); \coordinate (D) at (4,1); \coordinate (E) at (1,0);
\coordinate (Y) at (5,3); \coordinate (Z) at (7,2); \coordinate (X) at (6,1);
\draw[-{Stealth},thick] (A) -- (B);
\node[left] at (A) {$A$}; \node[above] at (B) {$B$};
\foreach \p in {C,D,E,Y,Z,X} \fill (\p) circle(1.5pt);
\node[left] at (C) {$C$}; \node[left] at (D) {$D$}; \node[above] at (E) {$E$};
\node[right] at (Y) {$Y$}; \node[left] at (Z) {$Z$}; \node[right] at (X) {$X$};
\end{tikzpicture}
\end{center}
\end{Exo}

\begin{Exo}{2}{\starfull{diffgreen}\starempty\starempty}{Repr\'esenter}
La figure ci-dessous repr\'esente sept hexagones r\'eguliers et num\'erot\'es.
\begin{center}
\begin{tikzpicture}[scale=0.22]
\draw (2,0)-- (4,0)-- (5,1.73)-- (4,3.46)-- (2,3.46)-- (1,1.73)-- cycle;
\draw (3,1.73) node[draw,circle]{1};
\draw (4,3.46)-- (5,5.2)-- (4,6.93)-- (2,6.93)-- (1,5.2)-- (2,3.46);
\draw (3,5.2) node[draw,circle]{2};
\draw (5,1.73) --(7,1.73)-- (8,3.46)-- (7,5.2)-- (5,5.2);
\draw (6,3.46) node[draw,circle]{3};
\draw (4,0)-- (5,-1.73)-- (7,-1.73)-- (8,0)--(7,1.73);
\draw (6,0) node[draw,circle]{4};
\draw (8,0)-- (10,0)-- (11,1.73)-- (10,3.46)-- (8,3.46);
\draw (9,1.73) node[draw,circle]{5};
\draw (10,3.46)-- (11,5.2)-- (10,6.93)-- (8,6.93)-- (7,5.2);
\draw (9,5.2) node[draw,circle]{6};
\draw (8,6.93)-- (7,8.66)-- (5,8.66)-- (4,6.93);
\draw (6,6.93) node[draw,circle]{7};
\draw (2,3.46) node[left] {$A$};
\draw (5,8.66) node[left] {$B$};
\draw (8,6.93) node[above right] {$C$};
\draw (11,5.2) node[right] {$D$};
\draw (11,1.73) node[right] {$E$};
\end{tikzpicture}
\end{center}
D\'eterminer l'image~: de l'hexagone 1 dans la translation de vecteur $\Vj{AC}$~; de l'hexagone 4 dans la translation de vecteur $\Vj{AB}$~; de l'hexagone 7 dans la translation de vecteur $\Vj{DE}$.
\end{Exo}

\begin{Exo}{3}{\starfull{diffgreen}\starempty\starempty}{Repr\'esenter}
Sans justifier, donner l'image de chaque point par la translation de vecteur indiqu\'ee.
\begin{center}
\begin{tikzpicture}[scale=0.32]
\draw[thin,gray!35] (-1,-1.2) grid (11,4.6);
\coordinate (A) at (0,4); \coordinate (B) at (2,4); \coordinate (C) at (4,4); \coordinate (D) at (6,4); \coordinate (E) at (8,4); \coordinate (F) at (10,4);
\coordinate (G) at (-0.5,2); \coordinate (H) at (1.5,1.5); \coordinate (I) at (3.5,1.5); \coordinate (J) at (5.5,1.5); \coordinate (K) at (7.5,1.5); \coordinate (L) at (10.5,2);
\coordinate (M) at (0,-0.5); \coordinate (N) at (2,-0.5); \coordinate (O) at (4,-0.5); \coordinate (P) at (6,-0.5); \coordinate (Q) at (8,-0.5); \coordinate (R) at (10,-0.5);
\foreach \p in {A,B,C,D,E,F,G,H,I,J,K,L,M,N,O,P,Q,R} \fill (\p) circle(1.3pt);
\node[above] at (A) {$A$}; \node[above] at (B) {$B$}; \node[above] at (C) {$C$}; \node[above] at (D) {$D$}; \node[above] at (E) {$E$}; \node[above] at (F) {$F$};
\node[left] at (G) {$G$}; \node[below] at (H) {$H$}; \node[below] at (I) {$I$}; \node[below] at (J) {$J$}; \node[below] at (K) {$K$}; \node[right] at (L) {$L$};
\node[below] at (M) {$M$}; \node[below] at (N) {$N$}; \node[below] at (O) {$O$}; \node[below] at (P) {$P$}; \node[below] at (Q) {$Q$}; \node[below] at (R) {$R$};
\end{tikzpicture}
\end{center}
\begin{enumerate}
\item Image du point $I$ par la translation de vecteur $\Vj{AD}$.
\item Image du point $J$ par la translation de vecteur $\Vj{IG}$.
\item Image du point $R$ par la translation de vecteur $\Vj{KJ}$.
\end{enumerate}
\end{Exo}

\begin{Exo}{4}{\starfull{diffred}\starfull{diffred}\starfull{diffred}}{Raisonner}
Soit $\vec u$ un vecteur non nul, et $A$ un point. On note $A'$ l'image de $A$ par la translation de vecteur $\vec u$.
\begin{enumerate}
\item On consid\`ere un point $B$ tel que $\Vj{AB}$ ait la m\^eme direction et le m\^eme sens que $\vec u$, mais une norme diff\'erente de celle de $\vec u$. Le point $B$ peut-il \^etre \'egal \`a $A'$~? Justifier.
\item M\^eme question si $\Vj{AB}$ a la m\^eme direction que $\vec u$ mais le sens contraire.
\item \'Enoncer, \`a l'aide des trois caract\'eristiques d'un vecteur, une condition n\'ecessaire et suffisante pour que $B=A'$.
\end{enumerate}
\end{Exo}

\Objectif{2}{\'Egalit\'e de deux vecteurs, vecteur nul}

\begin{Exo}{5}{\starfull{diffgreen}\starempty\starempty}{Raisonner}
Sur la figure suivante, expliquer, en utilisant les termes direction, sens ou norme, pourquoi le vecteur $\Vj{AB}$ n'est \'egal \`a aucun des autres vecteurs repr\'esent\'es.
\begin{center}
\begin{tikzpicture}[scale=0.4]
\foreach \x in {0,...,7} \foreach \y in {0,...,4} \fill[gray!55] (\x,\y) circle(0.6pt);
\draw[-{Stealth},thick] (0,2) -- (1,3); \node[below] at (0,2) {$A$}; \node[above] at (1,3) {$B$};
\draw[-{Stealth},thick] (1,0) -- (3,0); \node[below] at (1,0) {$C$}; \node[below] at (3,0) {$D$};
\draw[-{Stealth},thick] (3,2) -- (5,4); \node[below] at (3,2) {$G$}; \node[above] at (5,4) {$H$};
\draw[-{Stealth},thick] (7,2) -- (6,1); \node[below] at (6,1) {$E$}; \node[above] at (7,2) {$F$};
\end{tikzpicture}
\end{center}
\end{Exo}

\begin{Exo}{6}{\starfull{diffgreen}\starempty\starempty}{Calculer}
$EFGH$ est un parall\'elogramme de centre $I$ repr\'esent\'e ci-contre.
\begin{center}
\begin{tikzpicture}[scale=0.55]
\coordinate (E) at (0,1.3); \coordinate (F) at (3,1.5); \coordinate (H) at (0.5,0); \coordinate (G) at (3.3,0.2);
\coordinate (I) at ($(E)!0.5!(G)$);
\draw[thick,teal!45] (E) -- (F) -- (G) -- (H) -- cycle; \draw[thick,teal!45] (E) -- (G); \draw[thick,teal!45] (F) -- (H);
\node[above] at (E) {$E$}; \node[above] at (F) {$F$}; \node[below] at (H) {$H$}; \node[below] at (G) {$G$}; \node[above] at (I) {$I$};
\end{tikzpicture}
\end{center}
\begin{enumerate}
\item Donner un vecteur \'egal \`a~: a. $\Vj{EF}$ \quad b. $\Vj{GH}$ \quad c. $\Vj{EH}$ \quad d. $\Vj{GF}$
\item Donner un vecteur oppos\'e au vecteur $\Vj{EF}$.
\item Compl\'eter par le point qui convient~: a. $\Vj{EI}=\Vj{I\dots}$ \quad b. $\Vj{HI}=\Vj{I\dots}$ \quad c. $\Vj{FI}=\Vj{I\dots}$
\item Compl\'eter~: $\Vj{EH}=-\Vj{\dots}$
\item Quel est le repr\'esentant d'origine $F$ du vecteur $\Vj{EH}$~? du vecteur $\Vj{IH}$~?
\end{enumerate}
\end{Exo}

\begin{Exo}{7}{\starfull{diffblue}\starfull{diffblue}\starempty}{Repr\'esenter}
\`A partir de la figure ci-dessous, citer un vecteur~: opp. \`a $\Vj{CD}$~; m\^eme direction et sens que $\Vj{AC}$~; m\^eme direction que $\Vj{BC}$ mais sens contraire~; \'egal au vecteur $\Vj{BA}$.
\begin{center}
\begin{tikzpicture}[scale=0.32]
\draw[thin,gray!35] (-0.5,-0.5) grid (9.5,8.5);
\coordinate (A) at (3,4); \coordinate (B) at (1,1); \coordinate (C) at (3,1); \coordinate (D) at (4,2);
\foreach \p in {A,B,C,D} \fill (\p) circle(1.5pt);
\node[left] at (A) {$A$}; \node[below] at (B) {$B$}; \node[below] at (C) {$C$}; \node[below] at (D) {$D$};
\draw[-{Stealth},thick] (0,3) -- (1,3) node[midway,above]{$\vec u$};
\draw[-{Stealth},thick] (1,5) -- (3,8) node[midway,above left]{$\vec v$};
\draw[-{Stealth},thick] (4,5) -- (4,7) node[midway,left]{$\vec w$};
\draw[-{Stealth},thick] (5,6) -- (6,7) node[midway,above]{$\vec r$};
\draw[-{Stealth},thick] (8,8) -- (6,3) node[midway,above]{$\vec s$};
\draw[-{Stealth},thick] (8,8) -- (7,8) node[midway,above]{$\vec t$};
\draw[-{Stealth},thick] (9,2) -- (8,1) node[midway,above]{$\vec p$};
\draw[-{Stealth},thick] (9,5) -- (9,4) node[midway,left]{$\vec m$};
\end{tikzpicture}
\end{center}
\end{Exo}

\noindent{\color{red}\itshape Voir le devoir maison.}\par\vspace{4pt}

\begin{Exo}{8}{\starfull{diffblue}\starfull{diffblue}\starempty}{Raisonner}
Soient $ABCD$ et $ABEF$ deux parall\'elogrammes.
\begin{center}
\begin{tikzpicture}[scale=0.5]
\coordinate (A) at (0,0); \coordinate (B) at (2.2,0); \coordinate (D) at (0.6,1.5); \coordinate (C) at (2.8,1.5);
\coordinate (F) at (0.9,2.2); \coordinate (E) at (3.1,2.2);
\draw[green!45!black] (A)--(B)--(C)--(D)--cycle;
\draw[teal] (A)--(B)--(E)--(F)--cycle;
\node[below] at (A) {$A$}; \node[below] at (B) {$B$}; \node[left] at (D) {$D$}; \node[right] at (C) {$C$};
\node[left] at (F) {$F$}; \node[right] at (E) {$E$};
\end{tikzpicture}
\end{center}
\begin{enumerate}
\item Donner deux vecteurs \'egaux au vecteur $\Vj{AB}$.
\item En d\'eduire la nature du quadrilat\`ere $EFDC$.
\end{enumerate}
\end{Exo}

\begin{Exo}{9}{\starfull{diffblue}\starfull{diffblue}\starempty}{Raisonner}
Sur la figure ci-contre, $ABCD$ est un rectangle, et $AEBC$ et $ACDF$ sont des parall\'elogrammes.
\begin{center}
\begin{tikzpicture}[scale=0.5]
\coordinate (D) at (0,2); \coordinate (C) at (3,2); \coordinate (A) at (0,0); \coordinate (B) at (3,0);
\coordinate (F) at (-1.8,1); \coordinate (E) at (0,-1.8);
\draw[blue] (D)--(C)--(B)--(A)--cycle;
\draw[green!50!black] (F)--(A)--(C)--(D)--cycle;
\draw[red] (F)--(A)--(E)--(B);
\node[above] at (D) {$D$}; \node[above] at (C) {$C$}; \node[below] at (A) {$A$}; \node[below] at (B) {$B$};
\node[left] at (F) {$F$}; \node[below] at (E) {$E$};
\end{tikzpicture}
\end{center}
\begin{enumerate}
\item Montrer que $\Vj{EB}=\Vj{FD}$.
\item Montrer que le quadrilat\`ere $EBDF$ est un parall\'elogramme de centre $A$.
\end{enumerate}
\end{Exo}

\begin{Exo}{10}{\starfull{diffred}\starfull{diffred}\starfull{diffred}}{Raisonner}
Soient $A$, $B$, $C$ et $D$ quatre points du plan tels que $\Vj{AB}=\Vj{CD}$.
D\'emontrer que $\Vj{AC}=\Vj{BD}$.
\end{Exo}

\Objectif{3}{Repr\'esentants d'un vecteur}

\begin{Exo}{11}{\starfull{diffgreen}\starempty\starempty}{Repr\'esenter}
Sur la figure ci-dessous,
\begin{enumerate}
\item citer deux repr\'esentants du vecteur $\Vj{AB}$, autres que lui-m\^eme~;
\item indiquer si le vecteur $\Vj{EF}$ est aussi un repr\'esentant du vecteur $\Vj{AB}$, en justifiant.
\end{enumerate}
\begin{center}
\begin{tikzpicture}[scale=0.42]
\draw[thin,gray!35] (-0.5,-0.5) grid (7.5,6.5);
\coordinate (A) at (0,1); \coordinate (B) at (2,2);
\coordinate (C) at (3,4); \coordinate (D) at (5,5);
\coordinate (G) at (5,0); \coordinate (H) at (7,1);
\coordinate (E) at (1,4); \coordinate (F) at (2,6);
\draw[-{Stealth},thick,blue] (A) -- (B); \draw[-{Stealth},thick,blue] (C) -- (D); \draw[-{Stealth},thick,blue] (G) -- (H);
\draw[-{Stealth},thick,red] (E) -- (F);
\node[below] at (A) {$A$}; \node[right] at (B) {$B$};
\node[left] at (C) {$C$}; \node[right] at (D) {$D$};
\node[below] at (G) {$G$}; \node[right] at (H) {$H$};
\node[left] at (E) {$E$}; \node[above] at (F) {$F$};
\end{tikzpicture}
\end{center}
\end{Exo}

\begin{Exo}{12}{\starfull{diffred}\starfull{diffred}\starfull{diffred}}{Raisonner}
{\small \`A traiter une fois l'objectif 4 vu (relation de Chasles).}\\
Soit $\vec u$ un vecteur non nul et $A$ un point. On construit les points $B$, $C$ et $D$ tels que $\Vj{AB}=\Vj{BC}=\Vj{CD}=\vec u$. D'apr\`es la relation de Chasles, $\Vj{AC}=\Vj{AB}+\Vj{BC}$ et $\Vj{AD}=\Vj{AB}+\Vj{BC}+\Vj{CD}$.
\begin{enumerate}
\item Exprimer $\Vj{AC}$ et $\Vj{AD}$ en fonction de $\vec u$.
\item Les vecteurs $\Vj{AC}$ et $\Vj{AD}$ sont-ils des repr\'esentants de $\vec u$~? Justifier.
\end{enumerate}
\end{Exo}

\Objectif{4}{Somme de deux vecteurs et relation de Chasles}

\begin{Exo}{13}{\starfull{diffgreen}\starempty\starempty}{Repr\'esenter}
On consid\`ere les vecteurs $\vec u$ et $\vec v$ ci-dessous, ainsi que les points $B$, $V$ et $I$.
\begin{center}
\begin{tikzpicture}[scale=0.5]
\draw[thin,gray!35] (0,0) grid (9,6);
\coordinate (Ref) at (4.5,5.5);
\draw[-{Stealth},thick,blue] (Ref) -- ++(0,-1.8) node[midway,right]{$\vec u$};
\draw[-{Stealth},thick,red] (Ref) -- ++(1.8,-0.8) node[midway,above]{$\vec v$};
\coordinate (B) at (1,4.5); \node[left] at (B) {$+\,B$}; \fill (B) circle(1.3pt);
\coordinate (V) at (7.5,1); \node[below] at (V) {$+\,V$}; \fill (V) circle(1.3pt);
\coordinate (I) at (1,1); \node[below] at (I) {$+\,I$}; \fill (I) circle(1.3pt);
\end{tikzpicture}
\end{center}
Construire~: a. le point $P$ tel que $\Vj{BP}=\vec u+\vec v$~; b. le point $G$ tel que $\Vj{VG}=\vec u+\vec v$~; c. le point $K$ tel que $\Vj{IK}=\vec u+\vec v$.
\end{Exo}

\begin{Exo}{14}{\starfull{diffblue}\starfull{diffblue}\starempty}{Calculer}
\'Ecrire chacune des sommes suivantes \`a l'aide d'un seul vecteur, en utilisant la relation de Chasles.
\begin{enumerate}
\item[a.] $\Vj{TR}+\Vj{UT}+\Vj{RU}$
\item[b.] $\Vj{ST}-\Vj{VU}+\Vj{TU}$
\end{enumerate}
\end{Exo}

\begin{Exo}{15}{\starfull{diffblue}\starfull{diffblue}\starempty}{Calculer}
$ABCDEF$ est un hexagone r\'egulier de centre $O$.
\begin{center}
\begin{tikzpicture}[scale=0.55]
\foreach \i/\lab/\pos in {0/A/below right,60/B/above right,120/C/above,180/D/above left,240/E/below left,300/F/below}
\coordinate (\lab) at (\i:1.4);
\draw (A)--(B)--(C)--(D)--(E)--(F)--cycle;
\draw (A)--(D); \draw (B)--(E); \draw (C)--(F);
\node[below right] at (A) {$A$}; \node[above right] at (B) {$B$}; \node[above] at (C) {$C$};
\node[above left] at (D) {$D$}; \node[below left] at (E) {$E$}; \node[below] at (F) {$F$};
\node at (0,0) {$O$};
\end{tikzpicture}
\end{center}
Compl\'eter les \'egalit\'es suivantes en n'utilisant que des noms de points de la figure.
\begin{multicols}{2}
\begin{enumerate}
\item $\Vj{AB}+\Vj{CD}=\dots$
\item $\Vj{OB}+\Vj{OE}=\dots$
\item $\Vj{FO}+\Vj{DO}=\dots$
\item $\Vj{CA}-\Vj{DC}=\dots$
\item $\Vj{AF}+\Vj{CD}=\dots$
\item $\Vj{AB}-\Vj{OA}+\Vj{AF}=\dots$
\end{enumerate}
\end{multicols}
\vspace{1pt}
\end{Exo}

\begin{Exo}{16}{\starfull{diffred}\starfull{diffred}\starfull{diffred}}{Raisonner}
Soit $ABC$ un triangle et $I$ le milieu de $[BC]$.
\begin{enumerate}
\item Justifier que $\Vj{IB}+\Vj{IC}=\vec 0$.
\item En utilisant la relation de Chasles, d\'emontrer que $\Vj{AB}+\Vj{AC}=2\Vj{AI}$.
\end{enumerate}
\begin{center}
\begin{tikzpicture}[scale=0.55]
\coordinate (A) at (1,3); \coordinate (B) at (0,0); \coordinate (C) at (4,0); \coordinate (I) at (2,0);
\draw (A)--(B)--(C)--cycle; \draw[dashed] (A)--(I);
\fill (A) circle(1.3pt); \fill (B) circle(1.3pt); \fill (C) circle(1.3pt); \fill (I) circle(1.3pt);
\node[above] at (A) {$A$}; \node[below] at (B) {$B$}; \node[below] at (C) {$C$}; \node[below] at (I) {$I$};
\end{tikzpicture}
\end{center}
\end{Exo}

\Objectif{5}{Produit d'un vecteur par un nombre r\'eel et colin\'earit\'e de deux vecteurs}

\begin{Exo}{17}{\starfull{diffblue}\starfull{diffblue}\starempty}{Calculer}
Exprimer le vecteur $\Vj{AC}$ en fonction du vecteur $\Vj{AB}$ dans chacun des cas suivants~:
\begin{center}
\begin{tikzpicture}[scale=0.55]
\draw[-{Stealth}] (-2,0) -- (3,0);
\fill (0,0) circle(1.2pt) node[above]{$A$}; \fill (2,0) circle(1.2pt) node[above]{$B$}; \fill (-1.5,0) circle(1.2pt) node[above]{$C$};
\end{tikzpicture}
\hspace{1.2cm}
\begin{tikzpicture}[xscale=1.2,yscale=0.5]
\draw[-{Stealth}] (-1,0) -- (2.5,0);
\fill (0,0) circle(1.2pt) node[below]{$A$}; \fill (0.51,0) circle(1.2pt) node[above]{$B$}; \fill (2,0) circle(1.2pt) node[above]{$C$};
\end{tikzpicture}\\[10pt]
\begin{tikzpicture}[xscale=1.2,yscale=0.5]
\draw[-{Stealth}] (-2,0) -- (1.2,0);
\fill (0,0) circle(1.2pt) node[below]{$A$}; \fill (0.51,0) circle(1.2pt) node[above]{$B$}; \fill (-1.5,0) circle(1.2pt) node[above]{$C$};
\end{tikzpicture}
\hspace{1.2cm}
\begin{tikzpicture}[scale=0.55]
\draw[-{Stealth}] (-1,0) -- (2,0);
\fill (0,0) circle(1.2pt) node[above]{$A$}; \fill (1,0) circle(1.2pt) node[above]{$B$}; \fill (-0.5,0) circle(1.2pt) node[below]{$C$};
\end{tikzpicture}\\[10pt]
\begin{tikzpicture}[scale=0.45]
\draw[-{Stealth}] (-1,0) -- (4.5,0);
\fill (0,0) circle(1.2pt) node[above]{$A$}; \fill (3.5,0) circle(1.2pt) node[above]{$B$}; \fill (1,0) circle(1.2pt) node[above]{$C$};
\end{tikzpicture}
\hspace{1.2cm}
\begin{tikzpicture}[xscale=0.7,yscale=0.45]
\draw[-{Stealth}] (-1,0) -- (4,0);
\fill (0,0) circle(1.2pt) node[above]{$A$}; \fill (3,0) circle(1.2pt) node[above]{$B$}; \fill (0.5,0) circle(1.2pt) node[below]{$C$};
\end{tikzpicture}
\end{center}
\end{Exo}

\begin{Exo}{18}{\starfull{diffblue}\starfull{diffblue}\starempty}{Repr\'esenter}
Dans chacun des cas suivants, repr\'esenter le vecteur $\Vj{AC}$ tel que~:
\begin{center}
{\small a. $\Vj{AC}=-\dfrac32\Vj{AB}$ \quad b. $\Vj{AC}=\dfrac78\Vj{AB}$ \quad c. $\Vj{AC}=\dfrac14\Vj{AB}$ \quad d. $\Vj{AC}=-\dfrac12\Vj{AB}$ \quad e. $\Vj{AC}=\dfrac23\Vj{AB}$ \quad f. $\Vj{AC}=-\dfrac43\Vj{AB}$}
\end{center}
(un segment $[AB]$ diff\'erent est donn\'e pour chaque cas)
\end{Exo}

\begin{Exo}{19}{\starfull{diffblue}\starfull{diffblue}\starempty}{Calculer}
Soient deux points $A$ et $B$. D\'eterminer la valeur des r\'eels $m$, $n$, $p$ et $q$ tels que~:
\begin{multicols}{2}
\begin{enumerate}
\item[a.] $3\Vj{AB}+\dfrac12\Vj{AB}=m\Vj{AB}$
\item[b.] $2(3\Vj{BA})=n\Vj{AB}$
\item[c.] $2\Vj{AB}-\Vj{BA}=p\Vj{AB}$
\item[d.] $3\Vj{BA}-4\Vj{AB}=q\Vj{AB}$
\end{enumerate}
\end{multicols}
\vspace{1pt}
\end{Exo}

\Objectif{6}{D\'ecomposition d'un vecteur en fonction de deux vecteurs non colin\'eaires}

\begin{Exo}{20}{\starfull{diffgreen}\starempty\starempty}{Repr\'esenter}
Sur la figure ci-dessous, exprimer les vecteurs $\vec{w_1}$, $\vec{w_2}$ et $\vec{w_3}$ en fonction de $\vec u$ et $\vec v$.
\begin{center}
\begin{tikzpicture}[scale=0.5]
\draw[thin,gray!35] (-1.5,-1.5) grid (3.5,3.5);
\draw[-{Stealth},thick] (0,0) -- (1,0) node[midway,below]{$\vec u$};
\draw[-{Stealth},thick] (0,0) -- (0,1) node[midway,left]{$\vec v$};
\draw[-{Stealth},thick,blue] (0,0) -- (3,2) node[midway,above]{$\vec{w_1}$};
\draw[-{Stealth},thick,red] (0,0) -- (-1,3) node[midway,above]{$\vec{w_2}$};
\draw[-{Stealth},thick,teal] (0,0) -- (2,-1) node[midway,below right]{$\vec{w_3}$};
\end{tikzpicture}
\end{center}
\end{Exo}

\begin{Exo}{21}{\starfull{diffblue}\starfull{diffblue}\starempty}{Calculer}
$ABCD$ est un parall\'elogramme et $I$ est le milieu de $[CD]$. On pose $\vec u=\Vj{AB}$ et $\vec v=\Vj{AD}$.
Exprimer le vecteur $\Vj{AI}$ en fonction de $\vec u$ et $\vec v$.
\begin{center}
\begin{tikzpicture}[scale=0.6]
\coordinate (A) at (0,0); \coordinate (B) at (3,0); \coordinate (D) at (0.8,2); \coordinate (C) at (3.8,2);
\coordinate (I) at ($(D)!0.5!(C)$);
\draw (A)--(B)--(C)--(D)--cycle; \draw[dashed] (A)--(I);
\node[below] at (A) {$A$}; \node[below] at (B) {$B$}; \node[above] at (D) {$D$}; \node[above] at (C) {$C$};
\node[above] at (I) {$I$};
\end{tikzpicture}
\end{center}
\end{Exo}

\begin{Exo}{22}{\starfull{diffred}\starfull{diffred}\starfull{diffred}}{Calculer}
On consid\`ere de nouveau l'hexagone $ABCDEF$ (r\'egulier, centre $O$) de l'exercice 15. On note $\vec u=\Vj{OA}$ et $\vec v=\Vj{OB}$. Exprimer les vecteurs suivants \`a l'aide des vecteurs $\vec u$ et $\vec v$ uniquement~: $\Vj{CD}$~; $\Vj{CB}$~; $\Vj{DB}$~; $\Vj{AB}$~; $\Vj{DF}$.
\end{Exo}

\Objectif{7}{Caract\'erisation vectorielle du milieu d'un segment}

\begin{Exo}{23}{\starfull{diffgreen}\starempty\starempty}{Raisonner}
Soient $A$ et $B$ deux points distincts et $I$ le milieu de $[AB]$.
\begin{enumerate}
\item Que peut-on dire des vecteurs $\Vj{AI}$ et $\Vj{IB}$ (direction, sens, norme)~?
\item En d\'eduire que $\Vj{AI}=\Vj{IB}$.
\end{enumerate}
\begin{center}
\begin{tikzpicture}[scale=0.7]
\coordinate (A) at (0,0); \coordinate (I) at (2,0); \coordinate (B) at (4,0);
\draw[-{Stealth},thick] (A)--(I); \draw[-{Stealth},thick] (I)--(B);
\fill (A) circle(1.3pt); \fill (I) circle(1.3pt); \fill (B) circle(1.3pt);
\node[below] at (A) {$A$}; \node[below] at (I) {$I$}; \node[below] at (B) {$B$};
\end{tikzpicture}
\end{center}
\end{Exo}

\begin{Exo}{24}{\starfull{diffblue}\starfull{diffblue}\starempty}{Raisonner}
Soient $A$ et $B$ deux points, $I$ le milieu de $[AB]$, et $M$ un point quelconque du plan.
En utilisant la relation de Chasles et le fait que $\Vj{IA}+\Vj{IB}=\vec 0$, d\'emontrer que~:
$$\Vj{MA}+\Vj{MB}=2\Vj{MI}.$$
\end{Exo}

\noindent{\color{red}\itshape Voir le devoir maison.}\par\vspace{4pt}

\begin{Exo}{25}{\starfull{diffred}\starfull{diffred}\starfull{diffred}}{Raisonner}
On consid\`ere trois points $A$, $B$ et $C$. On nomme $A'$ le sym\'etrique de $A$ par rapport \`a $C$ et $B'$ le sym\'etrique de $B$ par rapport \`a $C$.
\begin{enumerate}
\item Montrer que $BAB'A'$ est un parall\'elogramme.
\item a. Que peut-on dire de $A'B'$ et $AB$~? \quad b. Quelle propri\'et\'e de la sym\'etrie centrale a-t-on d\'emontr\'ee~?
\end{enumerate}
\end{Exo}

\end{multicols}

\end{document}
